\label{sec:methodology}

\subsection{Measuring County Splits in \textsc{bisect} Plans}

A \textit{county split} occurs when a county's population is divided
between two or more congressional districts. County splits are measured
by the number of distinct (county, district) pairs in the final plan
minus the number of counties: if Texas has 254 counties and the plan
assigns them to 38 districts with 60 county splits, there are $254 + 60 = 314$
county-district pairs.

For Section~203 purposes, the relevant measure is whether a
\textit{Section~203-covered county} is split. Non-covered counties can be
split without Section~203 administrative consequences. The metric used here
is:

\[
\text{Covered County Splits} = \sum_{c \in \text{covered counties}} \mathbf{1}[\text{county } c \text{ spans} > 1 \text{ district}]
\]

\subsection{\textsc{bisect} Mode Comparison}

Two \textsc{bisect} operating modes are relevant for Section~203:

\textbf{Standard geographic mode} (\texttt{--weights-override geographic}):
Edge weights are proportional to the length of the shared boundary between
adjacent census tracts, scaled by population. This mode optimizes for
geographic compactness without specific consideration of county boundaries.
County splits occur when the bisection algorithm determines that splitting
a county produces a more balanced population partition.

\textbf{County-weight mode} (\texttt{--weights-override county}):
Edge weights on tract-to-tract boundaries that cross county lines are
reduced by a penalty factor, making the algorithm reluctant to cut
cross-county edges. This directly encodes a preference for county
integrity into the optimization. County splits still occur where
necessary for population balance, but the algorithm minimizes them.

The county-weight mode is particularly relevant for Section~203 analysis
because it specifically targets the administrative unit (the county)
at which Section~203 coverage is determined and at which election
administration is organized.

\subsection{Community Integrity Metric}

Beyond county splits, a finer-grained measure is needed for the analysis
of language minority communities that are smaller than counties. For the
Vietnamese community in Santa Clara County (California) and the Spanish-speaking
colonias in Hidalgo County (Texas), the relevant unit is not the county
but the census tract cluster that constitutes the community.

The community integrity metric used here is the fraction of language-minority
census tracts in a defined community that are assigned to the plurality district
(the district that receives the most tracts from the community):

\[
\text{Community Integrity} = \frac{\max_d |\{t \in C : d(t) = d\}|}{|C|}
\]

where $C$ is the set of census tracts in the community, $d(t)$ is the
district assigned to tract $t$, and $d$ ranges over all districts. A score
of 1.0 means the community is perfectly intact (all tracts in one district);
a score of 0.5 means the community is split evenly between two districts.

\subsection{Dispersed vs.\ Concentrated Language Communities}

Not all Section~203-covered language communities benefit equally from
\textsc{bisect}'s compactness optimization. Two distinct geographic
patterns are observed:

\textbf{Concentrated communities}: Urban language minority communities
(Vietnamese in San Jose's Little Saigon, Chinese in San Francisco's
Chinatown, Spanish-speaking colonias along the Texas-Mexico border)
are geographically compact. \textsc{bisect}'s geographic compactness
optimization tends to keep these communities intact because their tracts
form contiguous clusters that the algorithm naturally assigns to a single
partition.

\textbf{Dispersed communities}: Rural and semi-rural Spanish-speaking
populations in Texas (outside the border region), Navajo populations
spread across multiple Arizona counties, and similar dispersed communities
present a more difficult case. Because the community is not geographically
concentrated, the compactness optimization does not create natural pressure
to keep it intact. These communities may be better served by the
\textsc{bisect} VRA mode (\texttt{--weights-override vra-aligned}), which
explicitly weights minority community edges higher to maintain community
cohesion.
